%قبل از شروع مسئله ترجیحاً چیزی ننویسید. اما اگر خواستید بنویسید، در قالب لیستی (فقرات و شمارش) بنویسید.
\شروع‌مسئله[12]  
% صورت اصلی مسئله را در این‌جا بنویسید.
\\
شما عالم کل هستید و شنگدباو نادانی بیش نیست بنابراین از شما کلی سؤال در رابطه با داده‌ساختار ها دارد. به ازای هر مسئله زیر که شنگدباو پرسیده‌است داده‌ساختار مناسب پیشنهاد دهید و توضیحی مختصر در مورد نحوه کار داده ساختار مورد نظر برای مسئله مورد نظر ارائه دهید. نمره‌دهی با توجه به میزان بهینه بودن داده‌ساختار مورد نظر در بدترین حالت انجام می‌شود.
%اگر مسئله چند قسمت داشت به این صورت عمل کنید:
\شروع{زیرمسائل}
\زیرمسئله[4]
ما تعدادی گوی داریم که روی آن‌ها اعداد متمایز نوشته‌شده است.هر سری یا به ما گوی جدیدی داده می‌شود و یا اینکه یکی از گوی‌هاییکه گرفته ایم را باید پس بدهیم یا اینکه از ما پرسیده می‌شود که چند تا از گوی های ما عددشان در بازه 
$[l, r]$
است.

\زیرمسئله[5] 
دنباله‌ای از اعداد 
$a_1, ..., a_n$
در ابتدا موجود است. در هر مرحله یا یک اندیس‌ را باید انتخاب کنیم و مقدار آن‌را با $x$‌ای که به ما داده شده‌باید تعویض کنیم و یا اینکه بازه 
$[l, r]$
به ما داده می‌شود و باید تعداد صفر‌های جلوی عدد مقدار زیر را خروجی دهیم:
$$\Pi_{l \le i \le r} a_i$$
مقدار بالا برابر ضرب تمامی مقادیر در بازه 
$[l,r]$
است.
\زیرمسئله[3] 
فرض کنید در ابتدای کار یکسری کلمات انگلیسی را بلدیم و هر سری یا کلمه انگلیسی جدیدی می‌آموزیم یا اینکه یک کلمه به ما گفته می‌شود و می‌خواهیم ببینیم آن‌را می‌دانیم یا نه (یعنی قبلا به ما آموخته شده یا در ابتدای کار می‌دانستیم)

%\زیرمسئله[0] % مسئلۀ ۱ت
\پایان{زیرمسائل}



%%%%%%%%%%%%%%%%%%%%%%%%%%%%%%%%%%%%%%%%%%%%%%%%%%%%%
% برخی دستورات مفید برای تمرین‌ها. %
%%%%%%%%%%%%%%%%%%%%%%%%%%%%%%%%%%%%%%%%%%%%%%%%%%%%%

% فرمول‌های نمونه:
%$f(x) = \Theta(x)$
%$T(x, y) \leq \log(x) + 2^y + \binom{2n}{n}$
% {\tt code\_function}


% لیست موارد بدون شماره‌گذاری:
%\شروع{فقرات}
%    \فقره مورد اول 
%        \شروع{فقرات}
%            \فقره مورد اول داخلی 
%                \شروع{فقرات}
%                    \فقره اولین مورد خیلی داخلی 
%                    \فقره دومین مورد خیلی داخلی 
%                \پایان{فقرات}
%            \فقره مورد دوم داخلی 
%        \پایان{فقرات}
%    \فقره مورد دوم 
%\پایان{فقرات}

% لیست موارد شماره‌گذاری شده:
%\شروع{شمارش}
%    \فقره یک مورد 
%    \فقره یک مورد دیگر 
%    \فقره باز هم یک مورد دیگر
%\پایان{شمارش}

% برای نوشتن یک فرمول چند خطی که یک شماره بخوره. برای حالت‌های دیگر به سایت زیر مراجعه کنید:
% https://www.overleaf.com/learn/latex/Aligning_equations_with_amsmath
%\begin{align} 
%    \begin{split}
%        (x+y)^3 &= (x+y)^2(x+y) \\
%                &= (x^2+2xy+y^2)(x+y) \\
%                &= (x^3+2x^2y+xy^2) + (x^2y+2xy^2+y^3) \\
%                &= x^3+3x^2y+3xy^2+y^3
%    \end{split}                                 
%\end{align}

% نوشتن ماتریس و آرایه دوبعدی
%\begin{align}
%    A = 
%    \begin{bmatrix}
%        A_{11} & A_{21} \\
%        A_{21} & A_{22}
%    \end{bmatrix}
%\end{align}

% اضافه کردن تصویر با شرح. برای اضافه کردن تصویر بدون شرح، فقط دستور درج‌تصویر را استفاده کنید.
%\شروع{شکل}[h!]
%\تنظیم‌ازوسط
%  \درج‌تصویر[پهنا=0.5\پهنای‌متن]{img.jpg}
%\شرح{شرح تصویر}
%\برچسب{تصویر:۱}
%\پایان{شکل}
%\پایان{زیرمسائل}

%کد زیر را ببینید.
%%\lstinputlisting{source_filename.py}
%\begin{latin}
%\begin{lstlisting}
%# Import the modules
%import sys
%import random
%
%ans = True
%
%while ans:
%question = raw_input("Ask the magic 8 ball a question: (press enter to quit) ")
%
%answers = random.randint(1,8)
%
%if question == "":
%sys.exit()
%
%elif answers == 1:
%print "It is certain"
%
%elif answers == 2:
%print "Outlook good"
%
%elif answers == 3:
%print "You may rely on it"
%
%elif answers == 4:
%print "Ask again later"
%
%elif answers == 5:
%print "Concentrate and ask again"
%
%elif answers == 6:
%print "Reply hazy, try again"
%
%elif answers == 7:
%print "My reply is no"
%
%elif answers == 8:
%print "My sources say no"
%\end{lstlisting}
%\end{latin}
